%%%%%%%%%%%%%%%%%%%%%%%%%%%%%%%%%%%%%%%%%
% Journal Article
% LaTeX Template
% Version 1.0 (25/8/12)
%
% This template has been downloaded from:
% http://www.LaTeXTemplates.com
%
% Original author:
% Frits Wenneker (http://www.howtotex.com)
%
% License:
% CC BY-NC-SA 3.0 (http://creativecommons.org/licenses/by-nc-sa/3.0/)
%
%%%%%%%%%%%%%%%%%%%%%%%%%%%%%%%%%%%%%%%%%

%----------------------------------------------------------------------------------------
%	PACKAGES AND OTHER DOCUMENT CONFIGURATIONS
%----------------------------------------------------------------------------------------

\documentclass[twoside]{article}

\usepackage{lipsum} % Package to generate dummy text throughout this template

\usepackage[sc]{mathpazo} % Use the Palatino font
\usepackage[T1]{fontenc} % Use 8-bit encoding that has 256 glyphs
\linespread{1.05} % Line spacing - Palatino needs more space between lines
\usepackage{microtype} % Slightly tweak font spacing for aesthetics

\usepackage[hmarginratio=1:1,top=32mm,columnsep=20pt]{geometry} % Document margins
\usepackage{multicol} % Used for the two-column layout of the document
\usepackage{hyperref} % For hyperlinks in the PDF

\usepackage[hang, small,labelfont=bf,up,textfont=it,up]{caption} % Custom captions under/above floats in tables or figures
\usepackage{booktabs} % Horizontal rules in tables
\usepackage{float} % Required for tables and figures in the multi-column environment - they need to be placed in specific locations with the [H] (e.g. \begin{table}[H])

\usepackage{lettrine} % The lettrine is the first enlarged letter at the beginning of the text
\usepackage{paralist} % Used for the compactitem environment which makes bullet points with less space between them

\usepackage{abstract} % Allows abstract customization
\renewcommand{\abstractnamefont}{\normalfont\bfseries} % Set the "Abstract" text to bold
\renewcommand{\abstracttextfont}{\normalfont\small\itshape} % Set the abstract itself to small italic text

\usepackage{titlesec} % Allows customization of titles
\titleformat{\section}[block]{\large\scshape\centering{\Roman{section}.}}{}{1em}{} % Change the look of the section titles 

\usepackage{fancyhdr} % Headers and footers
\pagestyle{fancy} % All pages have headers and footers
\fancyhead{} % Blank out the default header
\fancyfoot{} % Blank out the default footer
\fancyhead[C]{ $\bullet$ November 2012 $\bullet$ } % Custom header text
\fancyfoot[RO,LE]{\thepage} % Custom footer text

%----------------------------------------------------------------------------------------
%	TITLE SECTION
%----------------------------------------------------------------------------------------

\title{\vspace{-15mm}\fontsize{20pt}{10pt}\selectfont\textbf{3-Level Geometric Modeler: CSG, Subdivision Trees and Boundary Representation}} % Article title

\author{
\large
\textsc{Kai Cao}\iffalse \thanks{Thanks to Chee Yap for his help}\fi \\[2mm] % Your name
\normalsize New York University \\ % Your institution
\normalsize \href{mailto:kcao@cs.nyu.edu}{kcao@cs.nyu.edu} % Your email address
\vspace{-5mm}
}
\date{}

%----------------------------------------------------------------------------------------

\begin{document}

\maketitle % Insert title

\thispagestyle{fancy} % All pages have headers and footers

%----------------------------------------------------------------------------------------
%	ABSTRACT
%----------------------------------------------------------------------------------------


%----------------------------------------------------------------------------------------
%	ARTICLE CONTENTS
%----------------------------------------------------------------------------------------

\begin{multicols}{2} % Two-column layout throughout the main article text

\section{Introduction}

Solid modeling is distinguished from related areas of geometric modeling and computer graphics by its emphasis on physical fidelity. Together, the principles of geometric and solid modeling form the foundation of computer-aided design and in general support the creation, exchange, visualization, animation, interrogation, and annotation of digital models of physical objects. 

We will describe a polyhedron as three sets : $(V, E, F)$. $V$ represents the set of vertices of the polyhedron. $E$ represents the set of edges of the polyhedron. $F$ represents the set of faces of the polyhedron.

%------------------------------------------------

\section{Definitions}

 \subsection {Regularized set}
      A set $S\subset \mathbb{R}^3$ is a regularized set $\textrm{iff}$ $S=\overline{int (S)}$. All of our operations are on Regularized sets.
\subsection {Regularized set operation on regularized sets $A$ and $B$}

	\begin{itemize}
		\item $A \cup^* B = \overline{int (A \cup B)}$
		\item $A \cap^* B = \overline{int (A \cap B)}$
		\item $A -^* B = \overline{int (A - B)}$
		\item $A^{c*} B = \overline{int (A^c)}$
	\end{itemize} 
      \subsection {Solid}
      A solid $S$ is defined as a regularized subset of $\mathbb{R}^3$  ($S\subset\mathbb{R}^3$)\\
      

%------------------------------------------------

\section{CSG objects}

Just as a polyhedron can be described as the intersection of half spaces, a CSG (Constructive Solid Geometry) model describes a polyhedron as regularized set operation on primitives which are the simplest solid objects used for the representation. 

\subsection {Primitives}
Primitives in CSG models are all semi-algebraic sets as follows:
\begin{enumerate}
      \item Half Space: $H=\{(x,y,z)| a*x+b*y+c*z-d\ge0\}$
      \item Sphere : $S=\{(x,y,z)| (x-x_0)^2+ (y-y_0)^2+ (z-z_0)^2-r^2\le0\}$. 

      
\end{enumerate}

%------------------------------------------------

\section {Subdivision trees}
      \begin{enumerate}
            \item Box: B is a box $\mathrm{iff}$ $B \subset \mathbb{R}^3$, $B=\{ (x,y,z)|x_1\le x\le x_2, y_1\le y\le y_2, z_1\le z\le z_2\}$, which is a regularized subset of $\mathbb{R}^3$
            \item Elementary Box for a polyhedron $P$: 
                  \begin{enumerate} 
            \item Empty Box:
            a box $B$ is an Empty Box $\mathrm{iff}$ $B\cap P=\emptyset$.
            \item Full Box: 
            a box $B$ is a Full Box of Polyhedron $P$ $\mathrm{iff}$ $B\cap P = B$
            \item Vertex Box: 
            a box $B$ is a Vertex Box $\mathrm{iff}$  there exists only one vertex $v \in V \cap B$, s.t.\ $\forall e \in E$, if $e \cap B \not= \emptyset$, $v$ bounds $e$, and $\forall f \in F$, if $f \cap B \not= \emptyset$, $v$ bounds $F$. 
            \item Edge Box: 
            a box $B$ is an Edge Box $\mathrm{iff}$ $V \cap B = \emptyset$ and there exists only one edge $ e\in E$, $e\cap B \not= \emptyset, \forall f \in F$, if $f \cap B \not= \emptyset$, e bounds f. 
            \item Face Box:
            a box $B$ is a Face Box $\mathrm{iff}$ $\forall e \in E$ , $e \cap B = \emptyset$, and there exists one face $f \in F$, and $f \cap B \not= \emptyset$. 



                  \end{enumerate}
	   	 \item Subdivision tree :  A subdivision tree in $\mathbb{R}^3$ is an octree, in which each node points to a box and we divide each box pointed by a inner node into 8 children pointing to 8 subboxes. 
           	 \item Adequate subdivision tree for box $B$ and polyhedron $P$ $T = SD (B_0, P)$ in which SD is a function on a  box $B$ and a polyhedron $P$, the root of the tree  points to box $B_0$.  All leaves of $T$ point to an elementary box of the polyhedron $P$.

		\item Boolean operations on two adequate subdivision trees $T_1 = SD (B_1, P_1)$ and $T_2 = SD (B_2, P_2)$: 
		$T = M (BO, T_1, T_2)$ in which B is a Boolean operation  and  $T_1$ and $T_2$ are adequate subdivision trees. The result of the function $M$ is an adequate subdivision tree $T$.\\
		if $B_1 \not= B_2$, $T=Null$.\\
		if $B_1 = B_2 = B_0$, $T$ is an adequate subdivision tree whose root points to $B _0$.
		$BO$ are Boolean operations as union $(\cup)$ , intersection $(\cap)$  and difference $(-)$.\\
	 	
		We can also define:\\
		\begin{itemize}
			\item $T_1 \cup T_2 = M\ (\cup, T_1, T_2)$
			\item $T_1 \cap T_2 = M\ (\cap, T_1, T_2)$
			\item $T_1 - T_2 = M\ (-, T_1, T_2)$
		\end{itemize}
		
		
			
		
	  \item Subdivision Rules of a box $B_0$ corresponding to a polyhedron $P$:
		\begin{enumerate}
			\item 
			\item
		\end{enumerate}
	 	

% Operations on elementary boxes
	\item Boolean operation on an elementary box for two polyhedra $P_1$ and $P_2$\\
	Suppose box $B_e$ is an elementary box for both $P_1$ and $P_2$. The the result of the Boolean operation on an an elementary box for two polyhedra $P_1$ and $P_2$ is an adequate subdivision tree. 
		\begin{enumerate}
			\item 
		\end{enumerate}
			
      \end{enumerate}

%----------------------------------------------------------------------------------------

\end{multicols}

\end{document}
